%=============================================================
%  preamble.tex — packages, colours, custom commands
%  Included by main.tex via %=============================================================
%  preamble.tex — packages, colours, custom commands
%  Included by main.tex via %=============================================================
%  preamble.tex — packages, colours, custom commands
%  Included by main.tex via %=============================================================
%  preamble.tex — packages, colours, custom commands
%  Included by main.tex via \input{preamble}
%=============================================================

% ─── Core Packages ──────────────────────────────────────────
\usepackage[a4paper, margin=1in]{geometry}
\usepackage{parskip}
\usepackage{titlesec}
\usepackage{graphicx}
\usepackage{float}
\usepackage{caption}
\usepackage{xcolor}
\usepackage[T1]{fontenc}
\usepackage{lmodern}
\usepackage[expansion=false]{microtype}

\usepackage{hyperref}
\usepackage{fancyhdr}
\usepackage{enumitem}
\usepackage{booktabs}
\usepackage{amsmath}
\usepackage{ragged2e}

% ─── listings (syntax highlighting, no shell-escape needed) ─
\usepackage{listings}

% ─── Colour Palette ─────────────────────────────────────────
\definecolor{primary}{RGB}{31, 73, 125}       % Dark blue  — headings
\definecolor{accent}{RGB}{0, 112, 192}        % Mid blue   — labels
\definecolor{codebg}{RGB}{245, 245, 250}      % Code background
\definecolor{coderule}{RGB}{180, 180, 180}    % Code border

% Python syntax colours
\definecolor{kw}{RGB}{0,   0,   180}   % keywords      (blue)
\definecolor{str}{RGB}{163, 21,  21}   % strings       (red)
\definecolor{cmt}{RGB}{80,  160, 80}   % comments      (green)
\definecolor{num}{RGB}{9,   134, 88}   % numbers       (teal)
\definecolor{blt}{RGB}{148, 0,   211}  % builtins      (purple)

% ─── listings Python Style ───────────────────────────────────
\lstdefinestyle{python}{
    language         = Python,
    backgroundcolor  = \color{codebg},
    basicstyle       = \ttfamily\small,
    keywordstyle     = \color{kw}\bfseries,
    stringstyle      = \color{str},
    commentstyle     = \color{cmt}\itshape,
    numberstyle      = \tiny\color{gray},
    numbers          = left,
    numbersep        = 8pt,
    stepnumber       = 1,
    frame            = single,
    framesep         = 4pt,
    rulecolor        = \color{coderule},
    breaklines       = true,
    breakatwhitespace= true,
    showstringspaces = false,
    tabsize          = 4,
    captionpos       = b,
    % Extra keywords (common builtins / libraries)
    morekeywords     = {import, as, from, True, False, None,
                        print, len, range, type, int, str, float,
                        list, dict, tuple, set},
    % emph = highlight pandas/plt methods differently
    emph             = {pd, plt, df_student, np},
    emphstyle        = \color{blt},
    xleftmargin      = 14pt,
    xrightmargin     = 4pt,
    aboveskip        = 8pt,
    belowskip        = 4pt,
}

% Plain text output style (no colours, monospace box)
\lstdefinestyle{output}{
    backgroundcolor  = \color{white},
    basicstyle       = \ttfamily\small,
    frame            = single,
    rulecolor        = \color{coderule},
    numbers          = none,
    breaklines       = true,
    xleftmargin      = 14pt,
    xrightmargin     = 4pt,
    aboveskip        = 4pt,
    belowskip        = 4pt,
}

% Default style
\lstset{style=python}

% ─── Hyperlink Style ────────────────────────────────────────
\hypersetup{hidelinks}   % no colours, no boxes, nothing

% ─── Header / Footer ────────────────────────────────────────
\setlength{\headheight}{14pt}
\pagestyle{fancy}
\fancyhf{}
\fancyfoot[C]{\small\thepage}
\renewcommand{\headrulewidth}{0pt}

% ─── Section counter for correct subsection numbering ────────
\newcounter{expnum}
\setcounter{expnum}{0}

% ─── Section Styles ─────────────────────────────────────────
\titleformat{\section}
  {\large\bfseries\color{primary}}{}{0em}{}

% Subsections show as  1.1  1.2  2.1  2.2  etc.
\titleformat{\subsection}
  {\normalsize\bfseries\color{accent}}{\theexpnum.\arabic{subsection}.}{0.5em}{}

% ─── Spacing helper ─────────────────────────────────────────
\newcommand{\gap}{\vspace{8pt}}

% ─── Helper: Experiment header ───────────────────────────────
% Usage: \experiment{1}{Full title of experiment}
\newcommand{\experiment}[2]{%
    \newpage
    \stepcounter{expnum}%
    \setcounter{subsection}{0}%
    {\LARGE\bfseries Experiment #1}\\[4pt]%
    {\large\bfseries\color{primary}\parbox{\linewidth}{#2}}%
    \vspace{10pt}%
}

% ─── Question block (justified, no box) ─────────────────────
% Usage: \question{Write the question text here}
\newcommand{\question}[1]{%
    \vspace{4pt}%
    \begin{flushleft}%
        \parbox{\linewidth}{\justifying #1}%
    \end{flushleft}%
    \vspace{6pt}%
}
%=============================================================

% ─── Core Packages ──────────────────────────────────────────
\usepackage[a4paper, margin=1in]{geometry}
\usepackage{parskip}
\usepackage{titlesec}
\usepackage{graphicx}
\usepackage{float}
\usepackage{caption}
\usepackage{xcolor}
\usepackage[T1]{fontenc}
\usepackage{lmodern}
\usepackage[expansion=false]{microtype}

\usepackage{hyperref}
\usepackage{fancyhdr}
\usepackage{enumitem}
\usepackage{booktabs}
\usepackage{amsmath}
\usepackage{ragged2e}

% ─── listings (syntax highlighting, no shell-escape needed) ─
\usepackage{listings}

% ─── Colour Palette ─────────────────────────────────────────
\definecolor{primary}{RGB}{31, 73, 125}       % Dark blue  — headings
\definecolor{accent}{RGB}{0, 112, 192}        % Mid blue   — labels
\definecolor{codebg}{RGB}{245, 245, 250}      % Code background
\definecolor{coderule}{RGB}{180, 180, 180}    % Code border

% Python syntax colours
\definecolor{kw}{RGB}{0,   0,   180}   % keywords      (blue)
\definecolor{str}{RGB}{163, 21,  21}   % strings       (red)
\definecolor{cmt}{RGB}{80,  160, 80}   % comments      (green)
\definecolor{num}{RGB}{9,   134, 88}   % numbers       (teal)
\definecolor{blt}{RGB}{148, 0,   211}  % builtins      (purple)

% ─── listings Python Style ───────────────────────────────────
\lstdefinestyle{python}{
    language         = Python,
    backgroundcolor  = \color{codebg},
    basicstyle       = \ttfamily\small,
    keywordstyle     = \color{kw}\bfseries,
    stringstyle      = \color{str},
    commentstyle     = \color{cmt}\itshape,
    numberstyle      = \tiny\color{gray},
    numbers          = left,
    numbersep        = 8pt,
    stepnumber       = 1,
    frame            = single,
    framesep         = 4pt,
    rulecolor        = \color{coderule},
    breaklines       = true,
    breakatwhitespace= true,
    showstringspaces = false,
    tabsize          = 4,
    captionpos       = b,
    % Extra keywords (common builtins / libraries)
    morekeywords     = {import, as, from, True, False, None,
                        print, len, range, type, int, str, float,
                        list, dict, tuple, set},
    % emph = highlight pandas/plt methods differently
    emph             = {pd, plt, df_student, np},
    emphstyle        = \color{blt},
    xleftmargin      = 14pt,
    xrightmargin     = 4pt,
    aboveskip        = 8pt,
    belowskip        = 4pt,
}

% Plain text output style (no colours, monospace box)
\lstdefinestyle{output}{
    backgroundcolor  = \color{white},
    basicstyle       = \ttfamily\small,
    frame            = single,
    rulecolor        = \color{coderule},
    numbers          = none,
    breaklines       = true,
    xleftmargin      = 14pt,
    xrightmargin     = 4pt,
    aboveskip        = 4pt,
    belowskip        = 4pt,
}

% Default style
\lstset{style=python}

% ─── Hyperlink Style ────────────────────────────────────────
\hypersetup{hidelinks}   % no colours, no boxes, nothing

% ─── Header / Footer ────────────────────────────────────────
\setlength{\headheight}{14pt}
\pagestyle{fancy}
\fancyhf{}
\fancyfoot[C]{\small\thepage}
\renewcommand{\headrulewidth}{0pt}

% ─── Section counter for correct subsection numbering ────────
\newcounter{expnum}
\setcounter{expnum}{0}

% ─── Section Styles ─────────────────────────────────────────
\titleformat{\section}
  {\large\bfseries\color{primary}}{}{0em}{}

% Subsections show as  1.1  1.2  2.1  2.2  etc.
\titleformat{\subsection}
  {\normalsize\bfseries\color{accent}}{\theexpnum.\arabic{subsection}.}{0.5em}{}

% ─── Spacing helper ─────────────────────────────────────────
\newcommand{\gap}{\vspace{8pt}}

% ─── Helper: Experiment header ───────────────────────────────
% Usage: \experiment{1}{Full title of experiment}
\newcommand{\experiment}[2]{%
    \newpage
    \stepcounter{expnum}%
    \setcounter{subsection}{0}%
    {\LARGE\bfseries Experiment #1}\\[4pt]%
    {\large\bfseries\color{primary}\parbox{\linewidth}{#2}}%
    \vspace{10pt}%
}

% ─── Question block (justified, no box) ─────────────────────
% Usage: \question{Write the question text here}
\newcommand{\question}[1]{%
    \vspace{4pt}%
    \begin{flushleft}%
        \parbox{\linewidth}{\justifying #1}%
    \end{flushleft}%
    \vspace{6pt}%
}
%=============================================================

% ─── Core Packages ──────────────────────────────────────────
\usepackage[a4paper, margin=1in]{geometry}
\usepackage{parskip}
\usepackage{titlesec}
\usepackage{graphicx}
\usepackage{float}
\usepackage{caption}
\usepackage{xcolor}
\usepackage[T1]{fontenc}
\usepackage{lmodern}
\usepackage[expansion=false]{microtype}

\usepackage{hyperref}
\usepackage{fancyhdr}
\usepackage{enumitem}
\usepackage{booktabs}
\usepackage{amsmath}
\usepackage{ragged2e}

% ─── listings (syntax highlighting, no shell-escape needed) ─
\usepackage{listings}

% ─── Colour Palette ─────────────────────────────────────────
\definecolor{primary}{RGB}{31, 73, 125}       % Dark blue  — headings
\definecolor{accent}{RGB}{0, 112, 192}        % Mid blue   — labels
\definecolor{codebg}{RGB}{245, 245, 250}      % Code background
\definecolor{coderule}{RGB}{180, 180, 180}    % Code border

% Python syntax colours
\definecolor{kw}{RGB}{0,   0,   180}   % keywords      (blue)
\definecolor{str}{RGB}{163, 21,  21}   % strings       (red)
\definecolor{cmt}{RGB}{80,  160, 80}   % comments      (green)
\definecolor{num}{RGB}{9,   134, 88}   % numbers       (teal)
\definecolor{blt}{RGB}{148, 0,   211}  % builtins      (purple)

% ─── listings Python Style ───────────────────────────────────
\lstdefinestyle{python}{
    language         = Python,
    backgroundcolor  = \color{codebg},
    basicstyle       = \ttfamily\small,
    keywordstyle     = \color{kw}\bfseries,
    stringstyle      = \color{str},
    commentstyle     = \color{cmt}\itshape,
    numberstyle      = \tiny\color{gray},
    numbers          = left,
    numbersep        = 8pt,
    stepnumber       = 1,
    frame            = single,
    framesep         = 4pt,
    rulecolor        = \color{coderule},
    breaklines       = true,
    breakatwhitespace= true,
    showstringspaces = false,
    tabsize          = 4,
    captionpos       = b,
    % Extra keywords (common builtins / libraries)
    morekeywords     = {import, as, from, True, False, None,
                        print, len, range, type, int, str, float,
                        list, dict, tuple, set},
    % emph = highlight pandas/plt methods differently
    emph             = {pd, plt, df_student, np},
    emphstyle        = \color{blt},
    xleftmargin      = 14pt,
    xrightmargin     = 4pt,
    aboveskip        = 8pt,
    belowskip        = 4pt,
}

% Plain text output style (no colours, monospace box)
\lstdefinestyle{output}{
    backgroundcolor  = \color{white},
    basicstyle       = \ttfamily\small,
    frame            = single,
    rulecolor        = \color{coderule},
    numbers          = none,
    breaklines       = true,
    xleftmargin      = 14pt,
    xrightmargin     = 4pt,
    aboveskip        = 4pt,
    belowskip        = 4pt,
}

% Default style
\lstset{style=python}

% ─── Hyperlink Style ────────────────────────────────────────
\hypersetup{hidelinks}   % no colours, no boxes, nothing

% ─── Header / Footer ────────────────────────────────────────
\setlength{\headheight}{14pt}
\pagestyle{fancy}
\fancyhf{}
\fancyfoot[C]{\small\thepage}
\renewcommand{\headrulewidth}{0pt}

% ─── Section counter for correct subsection numbering ────────
\newcounter{expnum}
\setcounter{expnum}{0}

% ─── Section Styles ─────────────────────────────────────────
\titleformat{\section}
  {\large\bfseries\color{primary}}{}{0em}{}

% Subsections show as  1.1  1.2  2.1  2.2  etc.
\titleformat{\subsection}
  {\normalsize\bfseries\color{accent}}{\theexpnum.\arabic{subsection}.}{0.5em}{}

% ─── Spacing helper ─────────────────────────────────────────
\newcommand{\gap}{\vspace{8pt}}

% ─── Helper: Experiment header ───────────────────────────────
% Usage: \experiment{1}{Full title of experiment}
\newcommand{\experiment}[2]{%
    \newpage
    \stepcounter{expnum}%
    \setcounter{subsection}{0}%
    {\LARGE\bfseries Experiment #1}\\[4pt]%
    {\large\bfseries\color{primary}\parbox{\linewidth}{#2}}%
    \vspace{10pt}%
}

% ─── Question block (justified, no box) ─────────────────────
% Usage: \question{Write the question text here}
\newcommand{\question}[1]{%
    \vspace{4pt}%
    \begin{flushleft}%
        \parbox{\linewidth}{\justifying #1}%
    \end{flushleft}%
    \vspace{6pt}%
}
%=============================================================

% ─── Core Packages ──────────────────────────────────────────
\usepackage[a4paper, margin=1in]{geometry}
\usepackage{parskip}
\usepackage{titlesec}
\usepackage{graphicx}
\usepackage{float}
\usepackage{caption}
\usepackage{xcolor}
\usepackage[T1]{fontenc}
\usepackage{lmodern}
\usepackage[expansion=false]{microtype}

\usepackage{hyperref}
\usepackage{fancyhdr}
\usepackage{enumitem}
\usepackage{booktabs}
\usepackage{amsmath}
\usepackage{ragged2e}

% ─── listings (syntax highlighting, no shell-escape needed) ─
\usepackage{listings}

% ─── Colour Palette ─────────────────────────────────────────
\definecolor{primary}{RGB}{31, 73, 125}       % Dark blue  — headings
\definecolor{accent}{RGB}{0, 112, 192}        % Mid blue   — labels
\definecolor{codebg}{RGB}{245, 245, 250}      % Code background
\definecolor{coderule}{RGB}{180, 180, 180}    % Code border

% Python syntax colours
\definecolor{kw}{RGB}{0,   0,   180}   % keywords      (blue)
\definecolor{str}{RGB}{163, 21,  21}   % strings       (red)
\definecolor{cmt}{RGB}{80,  160, 80}   % comments      (green)
\definecolor{num}{RGB}{9,   134, 88}   % numbers       (teal)
\definecolor{blt}{RGB}{148, 0,   211}  % builtins      (purple)

% ─── listings Python Style ───────────────────────────────────
\lstdefinestyle{python}{
    language         = Python,
    backgroundcolor  = \color{codebg},
    basicstyle       = \ttfamily\small,
    keywordstyle     = \color{kw}\bfseries,
    stringstyle      = \color{str},
    commentstyle     = \color{cmt}\itshape,
    numberstyle      = \tiny\color{gray},
    numbers          = left,
    numbersep        = 8pt,
    stepnumber       = 1,
    frame            = single,
    framesep         = 4pt,
    rulecolor        = \color{coderule},
    breaklines       = true,
    breakatwhitespace= true,
    showstringspaces = false,
    tabsize          = 4,
    captionpos       = b,
    % Extra keywords (common builtins / libraries)
    morekeywords     = {import, as, from, True, False, None,
                        print, len, range, type, int, str, float,
                        list, dict, tuple, set},
    % emph = highlight pandas/plt methods differently
    emph             = {pd, plt, df_student, np},
    emphstyle        = \color{blt},
    xleftmargin      = 14pt,
    xrightmargin     = 4pt,
    aboveskip        = 8pt,
    belowskip        = 4pt,
}

% Plain text output style (no colours, monospace box)
\lstdefinestyle{output}{
    backgroundcolor  = \color{white},
    basicstyle       = \ttfamily\small,
    frame            = single,
    rulecolor        = \color{coderule},
    numbers          = none,
    breaklines       = true,
    xleftmargin      = 14pt,
    xrightmargin     = 4pt,
    aboveskip        = 4pt,
    belowskip        = 4pt,
}

% Default style
\lstset{style=python}

% ─── Hyperlink Style ────────────────────────────────────────
\hypersetup{hidelinks}   % no colours, no boxes, nothing

% ─── Header / Footer ────────────────────────────────────────
\setlength{\headheight}{14pt}
\pagestyle{fancy}
\fancyhf{}
\fancyfoot[C]{\small\thepage}
\renewcommand{\headrulewidth}{0pt}

% ─── Section counter for correct subsection numbering ────────
\newcounter{expnum}
\setcounter{expnum}{0}

% ─── Section Styles ─────────────────────────────────────────
\titleformat{\section}
  {\large\bfseries\color{primary}}{}{0em}{}

% Subsections show as  1.1  1.2  2.1  2.2  etc.
\titleformat{\subsection}
  {\normalsize\bfseries\color{accent}}{\theexpnum.\arabic{subsection}.}{0.5em}{}

% ─── Spacing helper ─────────────────────────────────────────
\newcommand{\gap}{\vspace{8pt}}

% ─── Helper: Experiment header ───────────────────────────────
% Usage: \experiment{1}{Full title of experiment}
\newcommand{\experiment}[2]{%
    \newpage
    \stepcounter{expnum}%
    \setcounter{subsection}{0}%
    {\LARGE\bfseries Experiment #1}\\[4pt]%
    {\large\bfseries\color{primary}\parbox{\linewidth}{#2}}%
    \vspace{10pt}%
}

% ─── Question block (justified, no box) ─────────────────────
% Usage: \question{Write the question text here}
\newcommand{\question}[1]{%
    \vspace{4pt}%
    \begin{flushleft}%
        \parbox{\linewidth}{\justifying #1}%
    \end{flushleft}%
    \vspace{6pt}%
}